\documentclass[12pt,a4paper]{ctexart}

% ---- 页面设置 ----
\usepackage[top=2.5cm, bottom=2.5cm, left=2.5cm, right=2.5cm]{geometry}
\usepackage{graphicx}
\usepackage{hyperref}
\usepackage{listings}
\usepackage{xcolor}
\usepackage{enumitem}
\usepackage{tikz}
\usepackage{float}
\usepackage{tabularx}
\usepackage{booktabs}
\usepackage{fancyhdr}
\usepackage{titlesec}

\hypersetup{
    colorlinks=true,
    linkcolor=blue,
    urlcolor=cyan,
    citecolor=green
}

\definecolor{codebg}{RGB}{245,245,245}
\definecolor{codeframe}{RGB}{200,200,200}

\lstset{
    basicstyle=\ttfamily\small,
    backgroundcolor=\color{codebg},
    frame=single,
    rulecolor=\color{codeframe},
    breaklines=true,
    numbers=left,
    numberstyle=\tiny\color{gray},
    tabsize=2,
    language=C
}

\pagestyle{fancy}
\fancyhf{}
\fancyhead[L]{\small 桌面智能交互终端技术报告}
\fancyhead[R]{\small STM32F407 + ASR PRO + ESP8266}
\fancyfoot[C]{\thepage}
\renewcommand{\headrulewidth}{0.5pt}

\title{\textbf{桌面智能交互终端技术报告}\\
       \large 基于 STM32F407 的多模态智能交互终端}
\author{}
\date{2026年8月}

\begin{document}

\maketitle
\thispagestyle{empty}

\begin{center}
\begin{tabular}{rl}
\textbf{主控芯片} & STM32F407VET6 (ARM Cortex-M4, 168MHz) \\
\textbf{语音识别} & 天问 ASR PRO (CI-03T 离线语音芯片) \\
\textbf{WiFi通信}   & ESP8266 (ESP-AT 固件, 802.11 b/g/n) \\
\textbf{显示屏}     & 2.8寸 TFT-LCD (240×320, ILI9341, 16位并口) \\
\textbf{触摸屏}     & GT9147 电容触摸 (I2C) \\
\textbf{操作系统}   & FreeRTOS (CMSIS-RTOS v2) \\
\textbf{图形库}     & LVGL v8.3 \\
\textbf{云端AI}     & 阿里云 DashScope (Qwen 3.8-Max) \\
\textbf{开发环境}   & STM32CubeIDE + GCC ARM \\
\end{tabular}
\end{center}

\vspace{1cm}

\begin{abstract}
\noindent 本项目设计并实现了一款桌面智能交互终端。系统以 STM32F407 为主控制器，集成 LVGL 图形界面、虚拟表情动画、ASR PRO 离线语音命令控制，以及通过 ESP8266 连接阿里云 DashScope 大模型实现的云端 AI 对话功能。系统采用 FreeRTOS 多任务架构，通过消息队列实现模块间解耦通信，支持触摸屏拼音输入与离线语音命令两种交互方式，实现了具身化的智能对话体验。
\end{abstract}

\newpage
\tableofcontents
\newpage

% ============================================================
\section{系统总体架构}
% ============================================================

\subsection{硬件架构}

系统硬件由三大核心模块组成，通过串口和并行总线互联：

\begin{figure}[H]
\centering
\begin{tikzpicture}[node distance=2cm, auto]
    \tikzstyle{block} = [rectangle, draw, fill=blue!10, text width=3.5cm, text centered, rounded corners, minimum height=2cm]
    \tikzstyle{arrow} = [thick,->,>=stealth]

    \node [block] (stm32) {STM32F407VET6\\主控制器\\168MHz Cortex-M4};
    \node [block, above right=0.5cm and 2cm of stm32] (asr) {ASR PRO\\离线语音芯片\\CI-03T};
    \node [block, below right=0.5cm and 2cm of stm32] (esp) {ESP8266\\WiFi模块\\ESP-AT固件};
    \node [block, below=2cm of stm32] (lcd) {TFT-LCD\\2.8'' 240×320\\ILI9341 + GT9147};

    \draw [arrow] (asr) -- node[above] {USART1\\115200} (stm32);
    \draw [arrow] (stm32) -- node[below] {USART3\\115200} (esp);
    \draw [arrow] (stm32) -- node[left] {FSMC 16位并口} (lcd);
    \draw [arrow] (stm32) -- node[above,sloped] {I2C1} (lcd);
\end{tikzpicture}
\caption{系统硬件拓扑}
\end{figure}

\subsection{软件架构}

系统基于 FreeRTOS 实现多任务并行调度，共 4 个核心任务：

\begin{table}[H]
\centering
\begin{tabularx}{\textwidth}{lXXl}
\toprule
\textbf{任务名} & \textbf{功能} & \textbf{优先级} & \textbf{栈大小} \\
\midrule
defaultTask & 系统初始化、外设自检 & Normal & 4KB \\
lvglTask & LVGL图形刷新、UI事件处理、表情动画、ASR命令处理 & AboveNormal & 16KB \\
esp8266Task & ESP8266 AT指令交互、WiFi连接、HTTP/HTTPS请求 & Normal & 4KB \\
asrTask & ASR PRO串口数据接收与解析 & Normal & 4KB \\
\bottomrule
\end{tabular}
\caption{FreeRTOS 任务分配}
\end{table}

任务间通过 4 个消息队列通信：

\begin{itemize}
    \item \texttt{asr\_queue}：asrTask → lvglTask，传递语音命令 ID
    \item \texttt{esp\_ai\_req\_queue}：lvglTask → esp8266Task，传递 AI 对话请求
    \item \texttt{esp\_ai\_queue}：esp8266Task → lvglTask，传递 AI 回复文本
    \item \texttt{esp\_rx\_queue}：esp8266Task → lvglTask，传递系统状态日志
\end{itemize}

% ============================================================
\section{用户界面设计}
% ============================================================

\subsection{页面结构}

系统包含 5 个 LVGL 页面，通过语音命令或触摸操作切换：

\begin{table}[H]
\centering
\begin{tabularx}{\textwidth}{llX}
\toprule
\textbf{页面} & \textbf{语音命令} & \textbf{功能描述} \\
\midrule
待机画面 (idle\_scr) & ``待机画面'' & 虚拟表情动画（眨眼、眼球游走） \\
主菜单 (main\_scr) & ``打开菜单'' & 4 个功能入口按钮（2×2 布局） \\
系统 (system\_scr)  & ``系统状态'' / ``设置'' & ESP8266 连接状态日志、硬件信息 \\
时钟 (clock\_scr)   & ``时钟'' & 大字体数字时钟，每秒刷新 \\
AI对话 (ai\_scr)     & ``打开对话'' & 云端AI对话界面、拼音输入法 \\
\bottomrule
\end{tabular}
\caption{UI 页面说明}
\end{table}

\subsection{拼音输入法}

在 AI 对话页面实现了完整的拼音输入法，支持 GB2312 一级字库共 3755 个汉字：

\begin{itemize}
    \item 按拼音音节分组，405 个拼音码表，每组按使用频率排序候选字
    \item 8 个候选字按钮，实时匹配当前输入拼音串
    \item 支持模糊匹配（输入部分拼音即可显示候选字）
    \item 虚拟键盘输入，点击候选项自动替换末尾拼音串
\end{itemize}

\subsection{机器人面部表情动画}

系统实现了 4 种表情状态的状态机：

\begin{table}[H]
\centering
\begin{tabularx}{\textwidth}{llX}
\toprule
\textbf{状态} & \textbf{触发条件} & \textbf{动画效果} \\
\midrule
ROBOT\_IDLE & 默认状态、1.8秒无事件自动恢复 & 嘴部细线(80×5)闭合，眨眼(2-5秒随机)，眼球游走 \\
ROBOT\_LISTENING & 收到语音命令 & 嘴部微张(40×12)，提示正在聆听 \\
ROBOT\_THINKING & 提交AI请求 & 嘴部圆形(36×36)，表示思考中 \\
ROBOT\_SPEAKING & 收到AI回复 & 嘴部交替张合(60×36/60×8, 180ms周期)，模拟说话 \\
\bottomrule
\end{tabular}
\caption{机器人表情状态}
\end{table}

表情动画的核心参数：

\begin{itemize}
    \item 眨眼：2--5 秒随机间隔，闭合 80ms，张开 120ms
    \item 眼球游走：1.2--2.8 秒随机间隔，移动范围 ±7px，过渡 300ms
    \item 说话：180ms 周期交替张合
    \item 自动恢复：无事件 1.8 秒后切回 IDLE
\end{itemize}

% ============================================================
\section{语音控制模块}
% ============================================================

\subsection{ASR PRO 离线语音识别}

ASR PRO 是天问（TWen51）平台的离线语音识别芯片，基于 CI-03T，支持约 200 条命令词，无需联网即可进行语音识别。

\subsubsection{硬件接线}

\begin{table}[H]
\centering
\begin{tabular}{lll}
\toprule
\textbf{ASR PRO} & \textbf{STM32F407} & \textbf{说明} \\
\midrule
TX (PA9)  & USART1 RX (PA10) & 串口交叉连接 \\
RX (PA10) & USART1 TX (PA9)  & 串口交叉连接 \\
GND       & GND              & 共地 \\
\bottomrule
\end{tabular}
\caption{ASR PRO 与 STM32 接线}
\end{table}

\subsubsection{通信协议}

ASR PRO 识别到命令词后，通过串口发送 ASCII 十进制数字 + 换行符：

\begin{lstlisting}
// ASR PRO 端代码（上位机 GUI 自动生成）
void ASR_CODE() {
  set_state_enter_wakeup(10000);
  switch (snid) {
   case 1: Serial.write('1'); break;  // 打开菜单
   case 2: Serial.write('2'); break;  // 待机画面
   case 3: Serial.write('3'); break;  // 系统状态
   case 4: Serial.write('4'); break;  // AI对话
   case 5: Serial.write('5'); break;  // 返回
   case 6: Serial.write('6'); break;  // 设置
  }
  Serial.write('\n');
}
\end{lstlisting}

\subsubsection{STM32 端串口接收与解析}

STM32 通过 USART1 中断 + DMA 接收 ASR PRO 数据，实现了超时判帧机制：

\begin{itemize}
    \item 每收到一个字节，记录时间戳
    \item 收到换行符（0x0A）立即截帧
    \item 若 20ms 内无新字节且无换行符，自动截帧（兼容无换行符的 ASR 模块）
    \item 解析时遍历缓冲区，将 ASCII 十进制串转为整数 cmd\_id
    \item 解析完毕调用 \texttt{ASR\_FlushRx()} 清空缓冲区，重置 HAL 接收指针
\end{itemize}

% ============================================================
\section{云端 AI 对话模块}
% ============================================================

\subsection{通信链路}

\begin{center}
STM32F407 $\xrightarrow{\text{USART3, 115200}}$ ESP8266 $\xrightarrow{\text{WiFi}}$ 路由器 $\xrightarrow{\text{Internet}}$ 阿里云 DashScope
\end{center}

ESP8266 运行 ESP-AT 固件，作为透传网卡，STM32 通过 AT 指令控制其完成 WiFi 连接和 HTTPS 请求。

\subsection{AT 指令序列}

\begin{table}[H]
\centering
\begin{tabularx}{\textwidth}{lX}
\toprule
\textbf{指令} & \textbf{作用} \\
\midrule
\texttt{AT}               & 自检，确认模块存活 \\
\texttt{ATE0}             & 关闭回显，避免数据污染 \\
\texttt{AT+CWJAP="ssid","pwd"} & 连接 WiFi 热点 \\
\texttt{AT+CIPSSLSIZE=4096}    & 调整 SSL 缓冲区大小 \\
\texttt{AT+CIPSSLCSNI="host"}  & 设置 SNI（Server Name Indication） \\
\texttt{AT+CIPSTART="SSL","host",443} & 建立 SSL 加密连接 \\
\texttt{AT+CIPSEND=<len>} & 申请发送 HTTP 请求 \\
$\cdots$ HTTP 请求体 $\cdots$ & 发送 JSON 格式的 API 请求 \\
\bottomrule
\end{tabular}
\caption{ESP8266 AT 指令序列}
\end{table}

\subsection{HTTPS 通信的关键坑点}

\begin{enumerate}
    \item \textbf{SNI 必须设置}：阿里云 DashScope 走 CDN，同一 IP 对应多个域名。TLS 握手时 ESP8266 必须通过 \texttt{AT+CIPSSLCSNI} 告知服务器目标域名，否则服务器无法选择合适的证书，直接关闭连接。
    \item \textbf{确认 vs 响应}：\texttt{AT+CIPSEND} 的 ``Recv N bytes'' 是模块确认收到数据的字节数，不是远端响应。远端响应以 \texttt{+IPD} 开头。
    \item \textbf{超时设置}：大模型有冷启动延迟，首次请求可能超过 20 秒，接收超时需设为 60 秒。
    \item \textbf{响应解析}：HTTP 响应包含头部和 JSON body，需从 \texttt{"content"} 字段中提取 AI 回复文本，并处理 \texttt{\textbackslash n}、\texttt{\textbackslash "} 等转义字符。
\end{enumerate}

\subsection{API 调用}

调用阿里云 DashScope 兼容 OpenAI 格式的 Chat Completions API：

\begin{lstlisting}
POST https://dashscope.aliyuncs.com/compatible-mode/v1/chat/completions
Authorization: Bearer <API_KEY>
Content-Type: application/json

{
  "model": "qwen3.8-max",
  "messages": [
    {"role": "user", "content": "你好"},
    {"role": "assistant", "content": "你好！有什么可以帮助你的？"},
    {"role": "user", "content": "今天天气怎么样"}
  ]
}
\end{lstlisting}

系统维护最多 20 轮对话历史，每次请求将完整历史打包发送，实现多轮对话上下文。

% ============================================================
\section{关键代码模块}
% ============================================================

\subsection{项目文件结构}

\begin{lstlisting}[language=]
ai_project/
|-- Core/
|   |-- Inc/                    # 头文件
|   |   |-- FreeRTOSConfig.h    # FreeRTOS 配置
|   |   |-- pinyin_data.h       # 拼音输入法码表 (3755汉字)
|   |   |-- usart.h             # 串口配置
|   |   |-- fsmc.h              # FSMC 并口配置
|   |-- Src/
|       |-- main.c              # 入口，初始化外设
|       |-- freertos.c          # 核心：4个任务 + UI + 表情 + AI对话
|       |-- stm32f4xx_it.c      # 中断服务函数
|       |-- usart.c             # 串口初始化
|-- HARDWARE/
|   |-- ASR_PRO/
|   |   |-- asr_pro.c/h         # ASR PRO 串口驱动
|   |   |-- asrpro.txt          # ASR PRO 端代码参考
|   |-- ESP8266/
|   |   |-- esp8266.c/h         # ESP8266 AT 指令驱动
|   |-- LCD/
|   |   |-- lcd.c/h             # ILI9341 LCD 驱动
|   |-- LVGL/
|   |   |-- lv_port_disp.c/h    # LVGL 显示接口
|   |   |-- lv_port_indev.c/h   # LVGL 触摸输入接口
|   |-- ROBOT_FACE/
|   |   |-- robot_face_anim.c/h # 机器人表情动画
|   |-- TOUCH/
|       |-- gt9147.c/h          # GT9147 触摸驱动
|       |-- touch.c/h           # 触摸抽象层
|-- Drivers/                    # STM32 HAL 库 + CMSIS
|-- DEBUG.md                    # ESP8266 调试笔记
\end{lstlisting}

\subsection{LVGL 显示接口}

\begin{lstlisting}
// lv_port_disp.c - 双缓冲，240×20 行刷新
static lv_color_t buf1[240 * 20];
static lv_color_t buf2[240 * 20];

static void disp_flush(lv_disp_drv_t *drv, const lv_area_t *area,
                       lv_color_t *color_p) {
    LCD_Color_Fill(area->x1, area->y1, area->x2, area->y2, (u16 *)color_p);
    lv_disp_flush_ready(drv);
}
\end{lstlisting}

\subsection{ASR PRO 串口帧解析}

\begin{lstlisting}
// asr_pro.c - 超时判帧 + ASCII 十进制解析
static int ASR_Parse(asr_msg_t *msg) {
    if (!asr.data_ready || asr.rx_len == 0) return 0;

    uint32_t id = 0;
    for (int i = 0; i < asr.rx_len; i++) {
        uint8_t c = asr.rx_buf[i];
        if (c >= '0' && c <= '9')
            id = id * 10 + (c - '0');
        else
            break;
    }

    if (id == 0 || id > 160) { ASR_FlushRx(); return 0; }
    msg->cmd_id = (uint8_t)id;
    ASR_FlushRx();  // 重置缓冲区，重新开始接收
    return 1;
}
\end{lstlisting}

\subsection{任务间消息队列通信}

\begin{lstlisting}
// freertos.c - lvglTask 主循环 (简化)
for (;;) {
    // 1. 读取 ESP8266 状态消息，更新状态页面
    while (osMessageQueueGet(esp_rx_queueHandle, &m, NULL, 0) == osOK) { ... }

    // 2. 读取 AI 回复，追加到对话历史
    if (osMessageQueueGet(esp_ai_queueHandle, ai_buf, NULL, 0) == osOK) {
        chat_add_message("assistant", ai_buf);
        robot_face_set_state(ROBOT_SPEAKING);  // 触发说话表情
    }

    // 3. 读取 ASR 语音命令，切换页面
    if (osMessageQueueGet(asr_queueHandle, &asr_msg, NULL, 0) == osOK) {
        robot_face_set_state(ROBOT_LISTENING);
        switch (asr_msg.cmd_id) {
            case ASR_CMD_IDLE:  lv_scr_load(idle_scr); break;
            case ASR_CMD_MAIN_MENU: lv_scr_load(main_scr); break;
            // ... 其他页面切换
        }
    }

    // 4. 表情自动恢复：1.8秒无事件 → IDLE
    robot_face_tick();
    lv_timer_handler();
    osDelay(5);
}
\end{lstlisting}

% ============================================================
\section{调试与问题解决}
% ============================================================

\subsection{ASR PRO 串口通信问题}

\textbf{问题}：ASR PRO 语音命令无法切换 STM32 屏幕页面。

\textbf{根因}：ASR\_Parse 解析完一帧数据后仅将 \texttt{data\_ready} 置 0，未重置 \texttt{rx\_len} 和 HAL 接收指针。导致：
\begin{itemize}
    \item \texttt{rx\_len} 持续累积增长
    \item \texttt{HAL\_UART\_Receive\_IT} 指针越跑越偏，新数据存不到 \texttt{buf[0]}
    \item 解析器始终读取 \texttt{buf[0]} 的旧数据，同一命令被无限重复解析
\end{itemize}

\textbf{解决}：在 \texttt{ASR\_Parse} 末尾调用 \texttt{ASR\_FlushRx()}，一次性重置所有状态：\texttt{rx\_len=0}、清空缓冲区、\texttt{HAL\_UART\_Receive\_IT} 重新指向 \texttt{buf[0]}。

\subsection{ESP8266 HTTPS 连接问题}

详见 \texttt{DEBUG.md}，核心问题与解决方案：
\begin{itemize}
    \item \textbf{SNI 缺失}：ESP8266 老版本固件不支持 SNI，导致 TLS 握手失败，服务器直接关闭连接。通过在 \texttt{AT+CIPSTART} 前调用 \texttt{AT+CIPSSLCSNI} 解决。
    \item \textbf{确认误判}：将 \texttt{AT+CIPSEND} 的确认信息误判为远端响应，导致解析逻辑错误。
    \item \textbf{冷启动超时}：大模型首次调用需加载模型，超时从 20 秒调整为 60 秒。
\end{itemize}

\subsection{机器人表情嘴巴不闭合}

\textbf{问题}：待机状态下嘴巴看上去一直张着。

\textbf{根因}：\texttt{MOUTH\_IDLE\_H = 20}，20 像素高度的椭圆在视觉上像张开的嘴。

\textbf{解决}：将 \texttt{MOUTH\_IDLE\_H} 从 20 改为 5，圆角从 10 改为 2，形成一条细线，视觉上明确表示闭嘴状态。

% ============================================================
\section{技术指标}
% ============================================================

\begin{table}[H]
\centering
\begin{tabularx}{\textwidth}{lX}
\toprule
\textbf{指标} & \textbf{参数} \\
\midrule
CPU 主频 & 168MHz, Cortex-M4 \\
RAM 使用 & 约 80KB / 192KB (含 LVGL 双缓冲 9.6KB) \\
Flash 使用 & 约 400KB / 512KB (含 3755 汉字拼音码表) \\
LCD 刷新率 & 约 30--50 FPS (LVGL 双缓冲) \\
ASR 识别延迟 & < 500ms (离线) \\
AI 对话延迟 & 3--20 秒 (取决于网络和模型负载) \\
串口波特率 & 115200 (USART1 接 ASR PRO, USART3 接 ESP8266) \\
WiFi 标准 & 802.11 b/g/n, 2.4GHz \\
\bottomrule
\end{tabular}
\caption{系统技术指标}
\end{table}

% ============================================================
\section{总结与展望}
% ============================================================

\subsection{项目成果}

本项目成功实现了一款集语音控制、虚拟表情动画、云端 AI 对话于一体的桌面智能交互终端。主要成果包括：

\begin{enumerate}
    \item 基于 FreeRTOS 的多任务系统架构，任务间通过消息队列解耦，模块化程度高
    \item 基于 LVGL 的图形界面，包含 5 个页面，支持拼音输入法（3755 汉字）
    \item 机器人面部表情动画系统，4 种状态自动切换，具有情感化表达
    \item ASR PRO 离线语音命令控制，支持 7 条页面切换命令
    \item ESP8266 通过 HTTPS 连接阿里云 DashScope，实现多轮 AI 对话
\end{enumerate}

\subsection{未来优化方向}

\begin{enumerate}
    \item \textbf{语音输入升级}：将 ASR PRO 扩展为全双工语音对话，或替换为 ESP32-S3 实现云端流式语音识别
    \item \textbf{表情丰富化}：增加更多表情状态（开心、难过、惊讶、困惑），添加眉毛/腮红等面部元素
    \item \textbf{传感器扩展}：集成温湿度、光照、人体感应等传感器，实现环境感知
    \item \textbf{物联网功能}：通过 WiFi 获取天气、新闻、时间等信息，扩展实用功能
    \item \textbf{OTA 升级}：实现通过 WiFi 的固件远程升级
\end{enumerate}

\end{document}